\chapter{Operational Amplifiers}
\label{ch:opamps}

\section{The Ideal Op-Amp}

An \textbf{operational amplifier (op-amp)} is a high-gain differential
voltage amplifier IC, so widely used as a building block that it is
treated as a component in its own right rather than analyzed
transistor-by-transistor. It has two inputs --- the \textbf{non-inverting
input} ($+$) and the \textbf{inverting input} ($-$) --- and one output,
and amplifies the \emph{difference} between its two inputs by an
enormous open-loop gain (typically $10^5$ or more).

\bookfig[0.55]{opamp_internal.png}{A simplified internal structure of an
operational amplifier: differential input stage, gain stage, and output
stage.}

\begin{keyidea}[The two golden rules of ideal op-amp analysis]
When an op-amp is used with \textbf{negative feedback} (a path from the
output back to the inverting input), two simplifying assumptions make
hand analysis of nearly any op-amp circuit tractable:
\begin{enumerate}
\item \textbf{No current flows into either input} (the inputs have
essentially infinite impedance).
\item \textbf{The output adjusts itself so that the voltage difference
between the two inputs is zero} (``virtual short''), provided the
feedback loop is intact and the output is not saturated.
\end{enumerate}
\end{keyidea}

\section{The Inverting Amplifier}

\begin{center}
\begin{circuitikz}
\draw (0,2) to[R,l=$R_1$,-*] (2,2) node[op amp,anchor=-,scale=1.3](opamp){};
\draw (opamp.+) to[short] ++(-0.6,0) node[ground]{};
\draw (opamp.out) -- ++(0.8,0) node[right](out){$U_{\text{out}}$};
\draw (2,2) to[R,l=$R_f$] (2,3.2) -| (out.center) node[near end,above]{};
\draw (0,2) node[left]{$U_{\text{in}}$};
\end{circuitikz}
\end{center}

With the input applied through $R_1$ to the inverting input and feedback
resistor $R_f$ from output back to that same node, the ``virtual short''
rule forces the inverting input to sit at 0\,V (since the non-inverting
input is grounded), so the same current flows through $R_1$ and $R_f$.
The result is the widely used \textbf{inverting amplifier} gain formula:
\[
\boxed{A_v = \frac{U_{\text{out}}}{U_{\text{in}}} = -\frac{R_f}{R_1}}
\]
The negative sign indicates $180\degree$ phase inversion. Gain is set
purely by the ratio of two resistors, independent (to a good
approximation) of the op-amp's own large and imprecisely specified
open-loop gain --- the central practical benefit of negative feedback.

\begin{examplebox}[Inverting amplifier gain]
$R_1 = 10\,\text{k}\Omega$, $R_f = 100\,\text{k}\Omega$. Find the gain,
and the output voltage for a $-0.2\,\text{V}$ input.

\emph{Solution.}
\[
A_v = -\frac{R_f}{R_1} = -\frac{100}{10} = -10, \qquad
U_{\text{out}} = A_v \times U_{\text{in}} = -10 \times (-0.2) = 2\,\text{V}.
\]
\end{examplebox}

\section{The Non-Inverting Amplifier}

Applying the input directly to the non-inverting input, with a feedback
network ($R_f$ and $R_1$ to ground) setting the inverting input to a
fraction of the output, gives:
\[
\boxed{A_v = 1 + \frac{R_f}{R_1}}
\]
Gain is always $\geq 1$ and the output is \emph{in phase} with the
input (no inversion). Setting $R_f = 0$ (or removing it and connecting
output directly to the inverting input) gives $A_v = 1$: the
\textbf{voltage follower} (buffer), used purely for its extremely high
input impedance and low output impedance, without adding voltage gain.

\section{Other Common Op-Amp Circuits}

\begin{itemize}
\item \textbf{Summing amplifier}: multiple input resistors into the
inverting node combine several input voltages with individually
weighted gains, output being the (inverted, scaled) sum --- used in
audio mixing and analog computation.
\item \textbf{Differential amplifier}: amplifies the difference between
two input signals while rejecting any voltage common to both
(\textbf{common-mode rejection}) --- useful for extracting a small
signal riding on a larger common noise voltage.
\item \textbf{Integrator / differentiator}: replacing a resistor with a
capacitor in the feedback or input path produces a circuit whose output
is proportional to the time-integral or time-derivative of the input,
used in filter, waveform-shaping, and control applications.
\item \textbf{Active filters}: op-amps combined with RC networks
implement precise low-pass, high-pass, and band-pass filters (audio CW
and SSB filtering in receivers is a direct amateur radio application)
without the bulky inductors that a purely passive filter would need at
audio frequencies.
\item \textbf{Comparator}: an op-amp used without negative feedback (or
with positive feedback for hysteresis) simply drives its output to one
supply rail or the other depending on which input is higher --- used for
squelch circuits, threshold detection, and simple ADC building blocks.
\end{itemize}

\section{Real Op-Amp Limitations}

\begin{itemize}
\item \textbf{Finite open-loop gain and bandwidth}: gain is not
literally infinite, and it falls off with frequency (the
\textbf{gain-bandwidth product}, GBW, is a key datasheet specification
--- a 1\,MHz GBW device configured for a gain of 100 will only be
usable up to roughly 10\,kHz).
\item \textbf{Slew rate}: a maximum rate of change the output can
follow, in V/$\mu$s --- limits large-signal high-frequency performance
even within the small-signal bandwidth.
\item \textbf{Input offset voltage and bias current}: small
imperfections that cause a nonzero output even with both inputs
grounded, of particular concern in high-precision or high-gain DC
applications.
\item \textbf{Output swing and current limits}: the output cannot
exceed the supply rails (often not quite reaching them either, unless
the device is specifically a ``rail-to-rail'' design) and can only
supply a limited current before distorting or current-limiting.
\end{itemize}

Standard general-purpose op-amps (with gain-bandwidth products in the
single-digit MHz) are used extensively in amateur radio audio stages ---
receiver audio filtering, AGC (automatic gain control) generation, CW
sidetone oscillators, and microphone amplifiers --- but are generally
unsuitable as direct RF amplifiers at HF and above without specialized
high-speed devices, for which dedicated RF amplifier ICs or discrete
transistor stages (Chapter~\ref{ch:semiconductors}) are typically used
instead.

\begin{quickreview}
\item Ideal op-amp analysis: no current into either input; feedback
forces the two inputs to the same voltage (virtual short).
\item Inverting amplifier: $A_v = -R_f/R_1$ (inverted, gain set by
resistor ratio).
\item Non-inverting amplifier: $A_v = 1+R_f/R_1$ (in phase, gain
$\geq1$); $R_f=0$ gives a unity-gain buffer.
\item Summing, differential, integrator/differentiator, active filter,
and comparator circuits all build on the same two golden rules.
\item Real op-amps have finite gain-bandwidth product, slew rate, offset
errors, and output swing/current limits.
\end{quickreview}

\begin{practiceproblems}
\item Design an inverting amplifier with gain $-20$ using
$R_1=1\,\text{k}\Omega$; find the required $R_f$.
\item A non-inverting amplifier uses $R_1=2\,\text{k}\Omega$,
$R_f=18\,\text{k}\Omega$. Find the gain and the output for a 0.1\,V
input.
\item An op-amp has a gain-bandwidth product of 3\,MHz. Estimate the
maximum usable frequency at a closed-loop gain of 50.
\item Explain, using the golden rules, why the inverting input of an
inverting amplifier is called a ``virtual ground'' even though it is not
physically connected to ground.
\item Explain why a voltage-follower (buffer) stage is useful even
though its voltage gain is only 1.
\end{practiceproblems}
